%%%%%%%%%%  scan idx 45  |  folio 32  %%%%%%%%%%
composantes de $\ch(x)$ en degré $>0$ soient données par des polynômes universels isobares à coefficients rationnels en les composantes de $x$, et que
\begin{equation}\tag{1.27}
\ch([1,1+x^1])=\exp x^1=\sum_{n\geq0}(x^1)^n/n!.
\end{equation}
On obtient alors aussitôt l'expression explicite de $\ch$~: soit $\eta$ l'endomorphisme additif de $\widehat{A\otimes\Q}$ défini pour $x^i$ de degré $i$ par
\begin{equation}\tag{1.28}
\eta(x^i)=1/(-1)^{i-1}(i-1)!\,x^i
\end{equation}
on a alors
\begin{equation}\tag{1.29}
\ch\left([n,1+\sum_{i\geq1}x^i]\right)=n+\eta\left(\log\left(1+\sum_{i\geq1}x^i\right)\right).
\end{equation}
Comme on peut «~inverser~» cette formule, du moins si $A$ n'a qu'un nombre fini de composantes homogènes, on voit que dans ce cas, on a un \emph{isomorphisme} $\ch$ \emph{de l'anneau} $\widetilde A\otimes\Q$ \emph{sur l'anneau} $\widehat{A\otimes\Q}$ (ce qui élucide complètement la structure de $\widetilde A\otimes\Q$).

Rappelons que d'après Hirzebruch [1], une série formelle $f\in\Q[[t]]$ définit par des formules polynômiales universelles un homomorphisme additif $1+\Ahp\longrightarrow 1+(\widehat{A\otimes\Q})^+$, noté $\mathcal C_f$~; lorsque $f(t)=t/(1-\exp(-t))$, on pose simplement $\mathcal C_f=\mathcal C$ ($\mathcal C$ est l'initiale de Todd). On étend la définition de $\mathcal C(x)$ à $x\in\widetilde A$. Ceci dit, on a le résultat suivant~:

\textbf{Proposition 1.3.} \emph{Soit $N\in\widetilde A$ de la forme $[q,1+N^1+\cdots+N^q]$} (on a $N^i=0$ pour $i>q$). \emph{Si l'on pose~:}
\[
\lambda_{-1}(N)=\sum_{i\geq0}(-1)^i\lambda^i(N)=\sum_{i=0}^{n}(-1)^i\lambda^i(N),
\]
\emph{on aura~:}
\begin{equation}\tag{1.30}
\ch(\lambda_{-1}(N))=(-1)^q N^q\mathcal C(\check N)^{-1},
\end{equation}
